\chapter{Coalgebraic Minimization}\label{ch:coalgebra}

A molecule of Mol is a deterministic Mealy machine: a finite state space, an initial
state, and a transition $\mathit{step}\colon S \times A \to B \times S$ reading one
input symbol and emitting one output symbol (\autoref{ch:mol}). Read coalgebraically,
such a transition is a \emph{coalgebra} for a polynomial functor on sets: states carry
structure instead of being generated by operations, and the behavior of a state: the
outputs it produces on every future input word: is the image of the unique map into
a \emph{final} coalgebra. Minimization is then categorical, not syntactic: the minimal
state space is the Nerode quotient of the reachable states by behavioral
indistinguishability \citep{nerode}, computed by Hopcroft's partition refinement
\citep{hopcroft}.

The chapter proves four theorems. \autoref{thm:38}: every finite-state molecule has a minimal state
space, unique up to isomorphism; \autoref{thm:39}: minimization preserves behavior. Both restate
classical results of sequential-machine theory \citep{moore,mealy} and carry the tag
[C]. \autoref{thm:40} bounds the minimal state spaces of composed and tensored molecules by the
product of the component state spaces; \autoref{thm:41} shows behaviorally equivalent molecules
have isomorphic minimal state spaces; both are original and carry the tag [N].

\paragraph{Molecules as coalgebras.}
Fix an input alphabet $A$ and an output set $B$; in Mol these are the types of
\autoref{ch:mol}. For a set $X$, write $X^A$ for the functions $A \to X$. State spaces
are finite, as in the dataplane setting of \autoref{ch:mol}.

\begin{definition}[The machine functor]\label{def:machfun}
The \emph{machine functor} is the polynomial functor $F_{A,B}\colon \mathbf{Set} \to
\mathbf{Set}$, $F_{A,B} (X) = B \times X^A$, acting on $h\colon X \to Y$ by
$(b,g) \mapsto (b, h \circ g)$. A \emph{coalgebra} for $F_{A,B}$ is a pair $(S,c)$ with
$c \colon S \to F_{A,B} (S)$; writing $c (s) = (\mathit{out} (s), \mathit{nxt} (s))$,
$\mathit{out} (s) \in B$ is the output carried by $s$ and $\mathit{nxt} (s) \in S^A$ its
successor map. A \emph{homomorphism} $h \colon (S,c) \to (T,d)$ is a function
$h \colon S \to T$ with $d \circ h = F_{A,B} (h) \circ c$.
\end{definition}

\begin{definition}[Molecules as coalgebras]\label{def:molcoal}
A molecule presentation $(S, \mathit{step})$, $\mathit{step}\colon S \times A \to B
\times S$, is a coalgebra for the machine functor with output object $B^A$: currying
the input letter out of $\mathit{step}$ gives
\[
\widehat{\mathit{step}} \colon S \to (B \times S)^A \cong B^A \times S^A
= F_{A,B^A} (S),
\]
with output component $\mathit{out} (s) = (a \mapsto \mathit{out} (s,a)) \in B^A$ and
successor component $\mathit{nxt} (s) = (a \mapsto \mathit{nxt} (s,a)) \in S^A$. A
molecule of Mol is an isomorphism class of such coalgebras: the state-space bijections
of \autoref{ch:mol} are exactly the isomorphisms of coalgebras. A \emph{presentation}
carries an initial state $s_0 \in S$; the behavior of the molecule is the behavior of
$s_0$.
\end{definition}

\paragraph{Behavior and Nerode equivalence.}
For a state $s$, the \emph{response} $r_s \colon A^* \to B^*$ is defined by recursion:
$r_s (\varepsilon) = \varepsilon$ and
$r_s (aw) = \mathit{out} (s,a) \cdot r_{\mathit{nxt} (s,a)} (w)$; so $r_s (w)$ is the output
word produced from state $s$ on input $w$, and the behavior of the molecule is
$r_{s_0}$.

\begin{definition}[Nerode equivalence]\label{def:nerode}
States $s, t \in S$ are \emph{Nerode equivalent}, $s \sim t$, iff $r_s = r_t$
\citep{nerode}.
\end{definition}

\begin{lemma}[One-step refinement]\label{lem:onestep}
For all states $s, t$: $s \sim t$ iff for every letter $a \in A$,
$\mathit{out} (s,a) = \mathit{out} (t,a)$ and $\mathit{nxt} (s,a) \sim \mathit{nxt} (t,a)$.
\end{lemma}

\begin{proof}
($\Rightarrow$) If $r_s = r_t$, then $\mathit{out} (s,a) = r_s (a) = r_t (a) =
\mathit{out} (t,a)$ for every $a$. For every $v \in A^*$, the recursion gives
$r_s (av) = \mathit{out} (s,a) \cdot r_{\mathit{nxt} (s,a)} (v)$ and the analogous identity
for $t$; the equal words $r_s (av)$ and $r_t (av)$ have equal first symbols, so their
tails agree: $r_{\mathit{nxt} (s,a)} (v) = r_{\mathit{nxt} (t,a)} (v)$ for all $v$, i.e.,
$\mathit{nxt} (s,a) \sim \mathit{nxt} (t,a)$.

($\Leftarrow$) Suppose outputs agree at every letter and successors are equivalent.
Induct on $|w|$. For $|w| = 1$, $r_s (a) = \mathit{out} (s,a) = \mathit{out} (t,a) = r_t (a)$.
For $w = av$ with $v \neq \varepsilon$, the induction hypothesis and equal first symbols
give $r_s (av) = \mathit{out} (s,a) \cdot r_{\mathit{nxt} (s,a)} (v) = \mathit{out} (t,a)
\cdot r_{\mathit{nxt} (t,a)} (v) = r_t (av)$. Hence $r_s = r_t$.
\end{proof}

\paragraph{Final coalgebra semantics.}
Fix an output object $C$ and write $\Omega = C^{A^*}$ for the set of all functions
$A^* \to C$.

\begin{lemma}[Final coalgebra]\label{lem:final}
For the machine functor $F_{A,C} (X) = C \times X^A$, the pair
$(\Omega, \mathit{step}_\Omega)$ with
$\mathit{step}_\Omega (f) = (f (\varepsilon), (f_a)_{a \in A})$, where
$f_a (w) = f (aw)$, is the \emph{final coalgebra}: for every coalgebra $(S,c)$ there is
exactly one homomorphism $\mathsf{beh} \colon S \to \Omega$.
\end{lemma}

\begin{proof}
Each $f_a$ is a function $A^* \to C$, so $(\Omega, \mathit{step}_\Omega)$ is a
coalgebra. For $(S,c)$, $c (s) = (\mathit{out} (s), \mathit{nxt} (s))$, define
$\mathsf{beh} (s)(\varepsilon) = \mathit{out} (s)$ and
$\mathsf{beh} (s)(aw) = \mathsf{beh} (\mathit{nxt} (s,a))(w)$. Then
\[
\mathit{step}_\Omega (\mathsf{beh} (s)) =
(\mathsf{beh} (s)(\varepsilon), (w \mapsto \mathsf{beh} (s)(aw))_a)
= (\mathit{out} (s), (\mathsf{beh} (\mathit{nxt} (s,a)))_a)
= F_{A,C} (\mathsf{beh})(c (s)),
\]
so $\mathsf{beh}$ is a homomorphism. Any homomorphism $h$ is forced by the defining
equation to satisfy $h (s)(\varepsilon) = \mathit{out} (s)$ and
$h (s)(aw) = h (\mathit{nxt} (s,a))(w)$, determining $h (s)(w)$ by induction on $|w|$;
hence $h = \mathsf{beh}$.
\end{proof}

For a molecule, $C = B^A$: by induction on $|w|$,
$\mathsf{beh} (s)(w)(a) = \mathit{out} (\mathit{nxt} (s,w), a)$, the output emitted when
the state reached after reading $w$ reads $a$ next. Response and semantic map determine
each other: the $i$-th symbol of $r_s (w)$ is $\mathsf{beh} (s)(w_1 \cdots
w_{i-1})(w_i)$, and $\mathsf{beh} (s)(w)(a)$ is the last symbol of $r_s (wa)$: so
$s \sim t$ iff $\mathsf{beh} (s) = \mathsf{beh} (t)$: the Nerode equivalence is the
kernel of the final map.

\paragraph{The minimal state space.}
Let $(S, s_0, \mathit{step})$ be a finite-state presentation. A state $s$ is
\emph{reachable} if $s = \mathit{nxt} (s_0, w)$ for some word $w$; the reachable part
$S_0$ is finite and realizes the same behavior $r_{s_0}$. The quotient
\[
S_{\min} = S_0/{\sim}
\]
carries a coalgebra structure: by Lemma~\ref{lem:onestep}, $\mathit{out} (s,a)$ and the
class of $\mathit{nxt} (s,a)$ depend only on the class of $s$, so
$\mathit{out} ([s], a) = \mathit{out} (s,a)$ and
$\mathit{nxt} ([s], a) = [\mathit{nxt} (s,a)]$ are well defined. The pointed machine
$(S_{\min}, [s_0], \mathit{step}_{\min})$ is the \emph{minimization} $M_{\min}$ of
$M = (S, s_0, \mathit{step})$, with state space $S_{\min} (M)$. In the final coalgebra
it is the image $\mathsf{beh} (S_0)$, a subcoalgebra of $\Omega$ (successors of
reachable states are reachable), and $[s] \mapsto \mathsf{beh} (s)$ is an isomorphism of
pointed coalgebras onto it.

\setcatalog{38}
\begin{theorem}[Minimal state spaces (\autoref{thm:38})]\label{thm:38}
\textbf{[C]} Every finite-state molecule has a minimal state space, unique up to
isomorphism (Nerode/Hopcroft) \citep{nerode,hopcroft}.
\end{theorem}

\begin{proof}
We show $M_{\min}$ realizes the behavior of $M$ with the fewest states, and that any
two minimal realizations are isomorphic.

\emph{Existence.} By well-definedness, $M_{\min}$ reaches $[\mathit{nxt} (s_0, w')]$
after a prefix $w'$ of $w$ and emits $\mathit{out} ([s_{w'}], a) = \mathit{out} (s_{w'}, a)$
exactly where $M$ emits that symbol; induction on $|w|$ gives $r_{[s_0]} (w) = r_{s_0} (w)$
for all $w$.

\emph{Minimality.} Let $(T, t_0, \mathit{step}')$ realize $r_{s_0}$, with reachable
part $T_0$, and write $r'$ for responses in $T$. Put $s_w = \mathit{nxt} (s_0, w)$ and
$t_w = \mathit{nxt}'(t_0, w)$. Both machines emit the same output on $wu$; splitting
each response after $|w|$ symbols, the tail is the response of the state reached by
$w$, so $r'_{t_w} (u) = r_{s_w} (u)$ for every $u$. The map $t_w \mapsto [s_w] \colon
T_0 \to S_{\min}$ is well defined and surjective (every state of $S_0$ is $s_w$ for
some $w$), so $|T| \geq |T_0| \geq |S_{\min}|$.

\emph{Uniqueness.} If $(T, t_0, \mathit{step}')$ is minimal, then $|T| = |S_{\min}|$
and the surjection is a bijection $\varphi \colon T \to S_{\min}$ with $T_0 = T$. If
$\varphi (t) = [s]$, then $\mathit{out}'(t,a)$, the first symbol of $r'_t (a) = r_s (a)$,
equals $\mathit{out} (s,a) = \mathit{out}_{\min} (\varphi (t), a)$; and
$r'_{\mathit{nxt}'(t,a)} = r_{\mathit{nxt} (s,a)}$ (both are the tail of the common
response after $a$), so $\varphi (\mathit{nxt}'(t,a)) = [\mathit{nxt} (s,a)] =
\mathit{nxt}_{\min} (\varphi (t), a)$. Thus $\varphi$ is a homomorphism of pointed
coalgebras and, being a bijection, an isomorphism, with $\varphi (t_0) = [s_0]$. Every
minimal realization is isomorphic to $M_{\min}$. This is the classical theorem of
Nerode \citep{nerode}: the quotient by indistinguishability, following Moore
\citep{moore}: made effective by Hopcroft's refinement \citep{hopcroft}.
\end{proof}

\setcatalog{39}
\begin{theorem}[Minimization preserves behavior (\autoref{thm:39})]\label{thm:39}
\textbf{[C]} Minimization preserves behavior (Nerode quotient equivalence)
\citep{nerode}.
\end{theorem}

\begin{proof}
The states reached by $M$ and $M_{\min}$ on a word $w$ are $s_w$ and $[s_w]$, and every
emitted symbol agrees by well-definedness, so $r_{[s_0]} = r_{s_0}$: the minimized
machine produces exactly the behavior of $M$. In the language of Mol's behavioral
equivalence (\autoref{thm:6}), the relation
\[
R = \{ (s_w, [s_w]) : w \in A^* \}
\]
contains the initial pair $(s_0, [s_0])$ and is a bisimulation: related states emit
equal outputs on every letter and their successors $(s_{wa}, [s_{wa}])$ are again
related. Hence $M \simeq M_{\min}$. This is Nerode's quotient equivalence: minimization
quotients the state space by behavioral indistinguishability, and the quotient machine
is equivalent to the original.
\end{proof}

\setcatalog{40}
\begin{theorem}[Compositional state bounds (\autoref{thm:40})]\label{thm:40}
\textbf{[N]} For molecules $f$ and $g$ with state spaces $S_f$ and $S_g$,
\[
|S_{\min} (g \circ f)| \leq |S_f| \cdot |S_g|
\qquad\text{and}\qquad
|S_{\min} (f \otimes g)| \leq |S_f| \cdot |S_g| .
\]
\end{theorem}

\begin{proof}
Let $f$ be presented by $(S_f, s_f, \delta_f)$ and $g$ by $(S_g, s_g, \delta_g)$.

\emph{Sequential.} The composite $g \circ f$ is realized by the product machine with
state space $S_f \times S_g$, initial state $(s_f, s_g)$, and transition
$\delta_h ((u,v), a) = (c, (u',v'))$, where $(b, u') = \delta_f (u, a)$ and
$(c, v') = \delta_g (v, b)$ (\autoref{thm:49}). By \autoref{thm:38}, the minimal
realization has no more states than any realization, so $|S_{\min} (g \circ f)| \leq
|S_f \times S_g| = |S_f| \cdot |S_g|$.

\emph{Parallel.} The tensor $f \otimes g$ processes the components of an input pair
$(a, b)$ independently and is realized by the machine with state space $S_f \times S_g$
and transition $\delta_\otimes ((u,v),(a,b)) = (c,d),(u',v'))$, where
$(c, u') = \delta_f (u,a)$ and $(d, v') = \delta_g (v,b)$. The same minimality argument
gives $|S_{\min} (f \otimes g)| \leq |S_f| \cdot |S_g|$.

The bounds hold for any presentations of $f$ and $g$; they are sharpest when $S_f$ and
$S_g$ are minimal, since minimization discards unreachable states of the product
machine.
\end{proof}

\setcatalog{41}
\begin{theorem}[Minimization respects behavioral equivalence (\autoref{thm:41})]\label{thm:41}
\textbf{[N]} If $M \simeq M'$, then $S_{\min} (M) \cong S_{\min} (M')$.
\end{theorem}

\begin{proof}
Let $M$ and $M'$ be presented by $(S, s_0, \mathit{step})$ and
$(S', s'_0, \mathit{step}')$. Since $M \simeq M'$ (\autoref{thm:6}), some bisimulation
relates $s_0$ to $s'_0$; related states emit equal outputs on every letter and have
related successors, so by induction on $|w|$, $r_{s_0} (w) = r_{s'_0} (w)$ for all $w$:
both molecules realize the same behavior. Their minimizations are minimal realizations
of one common behavior, and by the uniqueness clause of \autoref{thm:38} any two minimal
realizations of the same behavior are isomorphic. Hence $S_{\min} (M) \cong S_{\min} (M')$
as pointed coalgebras.
\end{proof}

\paragraph{Partition refinement and Hopcroft's algorithm.}
The Nerode classes of a finite presentation are computable by \emph{partition
refinement}: start with the partition of $S_0$ by the output map $s \mapsto
\mathit{out} (s)$, and repeatedly split any block whose members send some letter into
different blocks. By Lemma~\ref{lem:onestep}, the partition by $\sim$ is the unique
coarsest stable partition refining the output partition: it is stable because
equivalent states have equivalent successors, and every stable refinement has each
block contained in one $\sim$-class (induction on word length: symbols agree by the
output refinement, successor blocks by stability). Hopcroft's algorithm computes this
partition in $O (k\,n \log n)$ steps for $n$ states and $k$ input letters
\citep{hopcroft}.

\paragraph{The coalgebraic reading.}
Minimization is categorical: the minimal coalgebra is the image of the unique map into
the final coalgebra, the quotient by the kernel of the semantic map being the coarsest
quotient coalgebra with the same behavior
\citep{arxiv-coalgebra-min,arxiv-coalgebra-min-2}; in Mol it is realized concretely by
$S_{\min} (M) \cong \mathsf{beh} (S_0)$. \autoref{thm:38}--\autoref{thm:39} give the minimal-state reading of
\autoref{thm:31} (finite-state molecules realize exactly the regular behaviors
\citep{brzozowski}); \autoref{thm:40} bounds state growth under the constructions of
\autoref{ch:mol}; \autoref{thm:41} shows minimization is well defined on Mol itself, descending to
Mol/$\simeq$ (\autoref{thm:6}). In practice this justifies compiling every hot-path
molecule to its minimal state space before building dispatch tables: behavior is
unchanged (\autoref{thm:39}), the tables have at most the product of the component sizes (\autoref{thm:40}),
and the result does not depend on how the molecule was written (\autoref{thm:41}).
